% GENERATED FILE. Edit canonical lessons, then run npm run generate:books.
% canonical-source-hash: fnv1a64:27204ead1bd9d1a8
% canonical-lessons: TE-C370-aitihyam, TE-C370-gatha, TE-C370-kirti, TE-C370-koniyadu, TE-C370-dandincu

\chapter{Legends}
\label{ch:te-legends}
\hlchaptermodality{370}

\begin{chapteropening}
\textbf{By the end of this chapter:} \emph{I can say a legend, a tale, fame, to extol and to chastise.}

Complete the last lesson of chapter 370: I can say a legend, a tale, fame, to extol and to chastise.
\end{chapteropening}

\section[\texorpdfstring{\te{ఐతిహ్యం} (aitihyaṁ)}{ఐతిహ్యం (aitihyaṁ)}]{\te{ఐతిహ్యం} (aitihyaṁ) — a legend; lore}

\label{lesson:TE-C370-aitihyam}



\begin{quote}
Before the new one: say the Telugu for to demolish, then the Telugu for to build.
\end{quote}

\subsection*{You'll want to know: \te{ఐతిహ్యం}}
\textbf{\te{ఐతిహ్యం}} — \emph{aitihyaṁ} — \textquotedblleft{}a legend; lore\textquotedblright{}.

\begin{quote}
\te{ఆ} \te{కోట} \te{గురించి} \te{ఒక} \te{ఐతిహ్యం} \te{ఉంది}.
\end{quote}

\begin{quote}
\emph{ā kōṭa guriñci oka aitihyaṁ undi.}
\end{quote}

\begin{quote}
There is a legend about that fort.
\end{quote}

A story handed down by word of mouth, not written in history books.

\begin{cousinweb}
\textbf{\te{ఐతిహ్యం}} (\emph{aitihyaṁ}) is borrowed from Sanskrit \emph{aitihya}, \textquotedblleft{}traditional instruction, tradition\textquotedblright{} (Monier-Williams), from \emph{iti ha}, \textquotedblleft{}thus indeed\textquotedblright{}.
\end{cousinweb}

\begin{grammarlens}[title={What you've built}]
One more word for telling what happened.
\end{grammarlens}

\subsection*{Your turn}
\begin{itemize}
\raggedright
  \item \emph{Say it:} \emph{aitihyaṁ}
  \item \emph{Say it:} \emph{aitihyaṁ}, once more
  \item \emph{Say it:} \emph{ā kōṭa guriñci oka aitihyaṁ undi.} — There is a legend about that fort.
  \item \emph{Recall:} say the Telugu for to demolish, then the Telugu for to build, then say \emph{aitihyaṁ} again
\end{itemize}

\begin{tcolorbox}[breakable,colback=teal!4,colframe=teal!35!black,title={Before you move on}]
What does \te{ఐతిహ్యం} mean? (\textquotedblleft{}A legend; lore\textquotedblright{}.) Say it once more.
\end{tcolorbox}
\section[\texorpdfstring{\te{గాథ} (gātha)}{గాథ (gātha)}]{\te{గాథ} (gātha) — a tale; a saga}

\label{lesson:TE-C370-gatha}



\begin{quote}
Before the new one: say the Telugu for to build, then the Telugu for a legend.
\end{quote}

\subsection*{You'll want to know: \te{గాథ}}
\textbf{\te{గాథ}} — \emph{gātha} — \textquotedblleft{}a tale; a saga\textquotedblright{}.

\begin{quote}
\te{మా} \te{అమ్మమ్మ} \te{వీరుల} \te{గాథలు} \te{చెప్పేది}.
\end{quote}

\begin{quote}
\emph{mā ammamma vīrula gāthalu ceppēdi.}
\end{quote}

\begin{quote}
My grandmother used to tell tales of heroes.
\end{quote}

Plural: \te{గాథలు}, \emph{gāthalu}.

\begin{cousinweb}
\textbf{\te{గాథ}} (\emph{gātha}) is borrowed from Sanskrit \emph{gāthā}, \textquotedblleft{}a song\textquotedblright{} (CDIAL 4126); in Telugu it is a tale of deeds.
\end{cousinweb}

\begin{grammarlens}[title={What you've built}]
One more word for telling what happened.
\end{grammarlens}

\subsection*{Your turn}
\begin{itemize}
\raggedright
  \item \emph{Say it:} \emph{gātha}
  \item \emph{Say it:} \emph{gātha}, once more
  \item \emph{Say it:} \emph{mā ammamma vīrula gāthalu ceppēdi.} — My grandmother used to tell tales of heroes.
  \item \emph{Recall:} say the Telugu for to build, then the Telugu for a legend, then say \emph{gātha} again
\end{itemize}

\begin{tcolorbox}[breakable,colback=teal!4,colframe=teal!35!black,title={Before you move on}]
What does \te{గాథ} mean? (\textquotedblleft{}A tale; a saga\textquotedblright{}.) Say it once more.
\end{tcolorbox}
\section[\texorpdfstring{\te{కీర్తి} (kīrti)}{కీర్తి (kīrti)}]{\te{కీర్తి} (kīrti) — fame; glory}

\label{lesson:TE-C370-kirti}



\begin{quote}
Before the new one: say the Telugu for a legend, then the Telugu for a tale.
\end{quote}

\subsection*{You'll want to know: \te{కీర్తి}}
\textbf{\te{కీర్తి}} — \emph{kīrti} — \textquotedblleft{}fame; glory\textquotedblright{}.

\begin{quote}
\te{ఆ} \te{కవి} \te{కీర్తి} \te{దేశమంతా} \te{వ్యాపించింది}.
\end{quote}

\begin{quote}
\emph{ā kavi kīrti dēśamantā vyāpiñcindi.}
\end{quote}

\begin{quote}
That poet's fame spread across the whole country.
\end{quote}

\te{కీర్తి} \te{గడించు}, to win fame.

\begin{cousinweb}
\textbf{\te{కీర్తి}} (\emph{kīrti}) is borrowed from Sanskrit \emph{kīrti}, \textquotedblleft{}report\textquotedblright{}, and then \textquotedblleft{}fame\textquotedblright{} (CDIAL 3200).
\end{cousinweb}

\begin{grammarlens}[title={What you've built}]
One more word for telling what happened.
\end{grammarlens}

\subsection*{Your turn}
\begin{itemize}
\raggedright
  \item \emph{Say it:} \emph{kīrti}
  \item \emph{Say it:} \emph{kīrti}, once more
  \item \emph{Say it:} \emph{ā kavi kīrti dēśamantā vyāpiñcindi.} — That poet's fame spread across the whole country.
  \item \emph{Recall:} say the Telugu for a legend, then the Telugu for a tale, then say \emph{kīrti} again
\end{itemize}

\begin{tcolorbox}[breakable,colback=teal!4,colframe=teal!35!black,title={Before you move on}]
What does \te{కీర్తి} mean? (\textquotedblleft{}Fame; glory\textquotedblright{}.) Say it once more.
\end{tcolorbox}
\section[\texorpdfstring{\te{కొనియాడు} (koniyāḍu)}{కొనియాడు (koniyāḍu)}]{\te{కొనియాడు} (koniyāḍu) — to extol; to praise highly}

\label{lesson:TE-C370-koniyadu}



\begin{quote}
Before the new one: say the Telugu for a tale, then the Telugu for fame.
\end{quote}

\subsection*{You'll want to know: \te{కొనియాడు}}
\textbf{\te{కొనియాడు}} — \emph{koniyāḍu} — \textquotedblleft{}to extol; to praise highly\textquotedblright{}.

\begin{quote}
\te{అందరూ} \te{ఆ} \te{వీరుడి} \te{ధైర్యాన్ని} \te{కొనియాడారు}.
\end{quote}

\begin{quote}
\emph{andarū ā vīruḍi dhairyānni koniyāḍāru.}
\end{quote}

\begin{quote}
Everyone praised that hero's courage.
\end{quote}

Past (I …): \te{కొనియాడాను}, \emph{koniyāḍānu}.

A written-style word: the news uses it of speeches and tributes.

\begin{cousinweb}
\textbf{\te{కొనియాడు}} (\emph{koniyāḍu}) is native Dravidian: DEDR 2151 lists Telugu \emph{koniyāḍu}, \textquotedblleft{}to praise\textquotedblright{}, under \emph{konu}, \textquotedblleft{}to take, hold\textquotedblright{}.
\end{cousinweb}

\begin{grammarlens}[title={What you've built}]
One more word for telling what happened.
\end{grammarlens}

\subsection*{Your turn}
\begin{itemize}
\raggedright
  \item \emph{Say it:} \emph{koniyāḍu}
  \item \emph{Say it:} \emph{koniyāḍu}, once more
  \item \emph{Say it:} \emph{andarū ā vīruḍi dhairyānni koniyāḍāru.} — Everyone praised that hero's courage.
  \item \emph{Recall:} say the Telugu for a tale, then the Telugu for fame, then say \emph{koniyāḍu} again
\end{itemize}

\begin{tcolorbox}[breakable,colback=teal!4,colframe=teal!35!black,title={Before you move on}]
What does \te{కొనియాడు} mean? (\textquotedblleft{}To extol; to praise highly\textquotedblright{}.) Say it once more.
\end{tcolorbox}
\section[\texorpdfstring{\te{దండించు} (daṇḍiñcu)}{దండించు (daṇḍiñcu)}]{\te{దండించు} (daṇḍiñcu) — to chastise; to discipline}

\label{lesson:TE-C370-dandincu}



\begin{quote}
Before the new one: say the Telugu for fame, then the Telugu for to extol.
\end{quote}

\subsection*{You'll want to know: \te{దండించు}}
\textbf{\te{దండించు}} — \emph{daṇḍiñcu} — \textquotedblleft{}to chastise; to discipline\textquotedblright{}.

\begin{quote}
\te{తండ్రి} \te{అల్లరి} \te{చేసిన} \te{కొడుకును} \te{దండించాడు}.
\end{quote}

\begin{quote}
\emph{taṇḍri allari cēsina koḍukunu daṇḍiñcāḍu.}
\end{quote}

\begin{quote}
The father punished his naughty son.
\end{quote}

Past (I …): \te{దండించాను}, \emph{daṇḍiñcānu}.

Built on the Sanskrit word for a rod.

\begin{cousinweb}
\textbf{\te{దండించు}} (\emph{daṇḍiñcu}) is built on Sanskrit \emph{daṇḍayati}, \textquotedblleft{}chastises, fines\textquotedblright{} (CDIAL 6136), from \emph{daṇḍa}, \textquotedblleft{}a stick, a rod\textquotedblright{} (CDIAL 6128).
\end{cousinweb}

\begin{grammarlens}[title={What you've built}]
One more word for telling what happened.
\end{grammarlens}

\subsection*{Your turn}
\begin{itemize}
\raggedright
  \item \emph{Say it:} \emph{daṇḍiñcu}
  \item \emph{Say it:} \emph{daṇḍiñcu}, once more
  \item \emph{Say it:} \emph{taṇḍri allari cēsina koḍukunu daṇḍiñcāḍu.} — The father punished his naughty son.
  \item \emph{Recall:} say the Telugu for fame, then the Telugu for to extol, then say \emph{daṇḍiñcu} again
\end{itemize}

\begin{tcolorbox}[breakable,colback=teal!4,colframe=teal!35!black,title={Before you move on}]
What does \te{దండించు} mean? (\textquotedblleft{}To chastise; to discipline\textquotedblright{}.) Say it once more.
\end{tcolorbox}
